% Part: history
% Chapter: biographies 
% Section: kurt-goedel

\documentclass[../../../include/open-logic-section]{subfiles}

\begin{document}

\olfileid[tr]{his}{bio}{god}

\olsection{Kurt G\"odel}

\olphoto{goedel-kurt}{Kurt G{\"o}del}

Kurt G{\"o}del (\textsc{ger}-dle), 28 Nisan 1906'da Avusturya-Macaristan
İmparatorluğu'ndaki Br{\"u}nn'de (bugün Çek Cumhuriyeti'ndeki Brno) doğdu.
Meraklı ve zeki yaradılışı nedeniyle genç Kurtele ailesi sık sık ``Der
kleine Herr Warum'' (Küçük Bay Neden) derdi. İlkokuldan başlayarak derslerinde
çok başarılıydı; en yüksek notun altında not aldığı tek ders matematikti.
Sağlığının kötü olması yüzünden G{\"o}del sık sık okula gidemez ve beden
eğitiminden muaf tutulurdu. Çocukluğunda romatizmal ateş tanısı aldı. Tıbbi
değerlendirmeler aksini söylese de bunun kalbini kalıcı biçimde etkilediğine
ömrü boyunca inandı.

G{\"o}del 1924'te Viyana Üniversitesi'nde öğrenime başladı ve doktora
çalışmalarını 1929'da tamamladı. Önce fizik okumayı düşünüyordu; fakat kısmen
filozof Rudolf Carnap'ın etkisiyle ilgisi kısa sürede matematiğe, özellikle de
mantığa yöneldi. Hans Hahn'ın danışmanlığında yazdığı tezinde, özdeşlik içeren
birinci dereceden yüklem mantığının tamlık teoremini ispatladı
\citep{Godel1929}. Yalnızca bir yıl sonra en ünlü sonuçlarını---birinci ve
ikinci eksiklik teoremlerini---elde etti (\citealt{Godel1931}'de yayımlandı).
Viyana'da bulunduğu yıllarda G{\"o}del, Carnap'ın da içinde yer aldığı bilim
yönelimli filozoflar topluluğu Viyana Çevresi'yle yakından ilişkiliydi;
özellikle Carnap'ın çalışmaları G{\"o}del'in sonuçlarından etkilendi.

G\"odel 1938'de Adele Nimbursky ile evlendi. Anne babası bu evlilikten hoşnut
değildi: Adele hem ondan altı yaş büyüktü ve boşanmıştı hem de bir gece
kulübünde dansçı olarak çalışıyordu. Ancak toplumsal baskılar G{\"o}del'i
etkilemedi ve çift onun ölümüne kadar mutlu bir evlilik sürdürdü.

Nazi Almanyası 1938'de Avusturya'yı ilhak ettikten sonra G{\"o}del ile Adele
Amerika Birleşik Devletleri'ne göç etti; G{\"o}del burada New Jersey,
Princeton'daki İleri Araştırmalar Enstitüsü'nde bir görev üstlendi. İçe dönük
ve sıra dışı yaradılışına karşın Princeton'daki yılları iş birliği içinde ve
verimli geçti. Küme teorisi, felsefe ve fizik üzerine yazılar yayımladı.
Özellikle de enstitüdeki meslektaşı Albert Einstein'la çok yakın bir dostluk
kurdu.

G{\"o}del'in ruh sağlığı ilerleyen yıllarında kötüleşti. Eşinin 1977'de
hastaneye yatırılması, artık onun için yemek pişirememesi demekti. Yaşamı
boyunca ruh sağlığı sorunları yaşayan G{\"o}del paranoyaya yenik düştü.
Zehirlenmekten ölümcül derecede korktuğu için yemek yemeyi reddetti. 14 Ocak
1978'de Princeton'da açlıktan öldü.

\begin{reading}
G{\"o}del'in yaşamının kapsamlı bir biyografisi için \citet{Dawson1997}
kaynağına bakın. Başka biyografik yazılar ile G{\"o}del'in mantık ve felsefeye
katkıları üzerine denemeler için \citet{Wang1990}, \citet{Baaz2011},
\citet{Takeuti2003} ve \citet{Sigmund2007} kaynaklarına bakın.

G{\"o}del'in doktora tezi özgün Almancasıyla mevcuttur \citep{Godel1929}.
Eksiklik teoremlerinin özgün metni \citep{Godel1931}'dir. G\"odel'in
yayımlanmış ve yayımlanmamış bütün yazıları ile yazışmalarından bir seçki,
İngilizce \emph{Collected Papers} içinde mevcuttur
\citet{Godel1986,Godel1990}.

G{\"o}del'in eksiklik teoremlerinin ayrıntılı bir ele alınışı için
\citet{Smith2013} kaynağına bakın. G{\"o}del'in teoremleri üzerine biçimsel
olmayan, felsefi bir tartışma için Mark Linsenmayer'in podcast'ine bakın
\citep{Linsenmayer2014}.
\end{reading}

\end{document}
